% Standalone machine note.
\documentclass[11pt]{article}
\usepackage[T1]{fontenc}
\usepackage[utf8]{inputenc}
\usepackage{amsmath,amssymb,mathtools}
\usepackage[margin=1in]{geometry}
\usepackage{microtype}
\usepackage{xcolor}
\usepackage[colorlinks=true,linkcolor=blue!55!black,urlcolor=blue!55!black]{hyperref}
\allowdisplaybreaks

\title{Analytic Weyl Anomaly from Plumbing to the Abel--Jacobi Metric}
\author{Machine Note}
\date{29 August 2026}

\newcommand{\dd}{\mathrm d}
\newcommand{\Tr}{\operatorname{Tr}}
\newcommand{\cP}{\mathcal P}
\newcommand{\cA}{\mathcal A}
\newcommand{\pl}{\mathrm{pl}}
\newcommand{\can}{\mathrm{can}}

\begin{document}
\maketitle

\paragraph{Purpose.}
We want the Weyl-anomaly factor that converts a vacuum partition function
defined by CFT sewing in a plumbing frame to the metric induced from the flat
Abel--Jacobi torus.  The direct Polyakov formula is awkward because a plumbing
metric is naturally singular at the seams.  Instead, we evaluate one real
free-boson partition function in the two frames and take their quotient.  The
result determines the Weyl factor of every CFT by its central charge.

\section{Abel--Jacobi map and canonical metric}

Let $\Sigma_g$ be a compact genus-$g$ Riemann surface.  Choose a symplectic
basis $(A_I,B_I)$ and normalized holomorphic one-forms
\begin{equation}
 \oint_{A_J}\omega_I=\delta_{IJ},
 \qquad
 \oint_{B_J}\omega_I=\Omega_{IJ},
 \qquad
 Y\equiv\operatorname{Im}\Omega>0.
\end{equation}
The Abel--Jacobi map is
\begin{equation}
 \mathfrak A_{p_0}:\Sigma_g\longrightarrow
 J(\Sigma_g)=\mathbb C^g/(\mathbb Z^g+\Omega\mathbb Z^g),
 \qquad
 p\longmapsto
 \left(\int_{p_0}^{p}\omega_1,\ldots,
       \int_{p_0}^{p}\omega_g\right).
\end{equation}
Pulling back the flat Hermitian form $\dd Z^T Y^{-1}\dd\bar Z$ gives the
Bergman, or Abel--Jacobi, metric.  We use its unit-area normalization,
\begin{equation}
 \boxed{
 \mu_{\can}
 =\frac{i}{2g}\sum_{I,J=1}^{g}
   (Y^{-1})_{IJ}\,\omega_I\wedge\overline{\omega_J},
 \qquad
 \int_{\Sigma_g}\mu_{\can}=1 .}
 \label{eq:canonical-metric}
\end{equation}
The unrescaled restriction has area $g$.  This fixed rescaling changes only a
moduli-independent determinant constant at fixed genus.

For one noncompact real boson, after dividing out the constant target-space
zero mode, write
\begin{equation}
 Z_X[g]=\left(\frac{\det{}'\Delta_g}{\operatorname{Area}(g)}\right)^{-1/2}.
 \label{eq:scalar-definition}
\end{equation}
If $g'=e^{2\sigma}g$ is smooth, then
\begin{equation}
 \log\frac{Z_X[g']}{Z_X[g]}
 =\frac{1}{24\pi}\int_{\Sigma_g}
 \left(|\nabla\sigma|_g^2+R_g\sigma\right)\dd A_g.
 \label{eq:polyakov}
\end{equation}
For plumbing metrics, the quotient below is the regularized definition of
the right-hand side, including its seam contributions.

\section{Problem and quotient strategy}

For a plumbing channel $a$, define
\begin{equation}
 W_a^{\pl/\can}
 \equiv\frac{Z_X^{\pl,a}}{Z_X^{\can}},
 \qquad
 \cA_a\equiv\log W_a^{\pl/\can}.
 \label{eq:weyl-definition}
\end{equation}
Here the plumbing propagator is the raw $q^{L_0}$ propagator: no
$q^{-c/24}$ Casimir factor is hidden in $Z_X^{\pl,a}$.  Momentum states obey
$\langle p|p'\rangle=\delta(p-p')$ and completeness uses $\dd p$.

Since the Weyl anomaly is proportional to the ordinary Virasoro central
charge $c$, a vacuum partition function transforms as
\begin{equation}
 \boxed{
 Z_c^{\can}=Z_c^{\pl,a}\left(W_a^{\pl/\can}\right)^{-c}
 =Z_c^{\pl,a}e^{-c\cA_a}.}
 \label{eq:central-charge-power}
\end{equation}

\section{Genus one}

For $\Sigma_1=\mathbb C/(\mathbb Z+\tau\mathbb Z)$, the Abel--Jacobi map is
the identity and
\begin{equation}
 \dd s_{\can}^2=\frac{|\dd z|^2}{\tau_2},
 \qquad \tau_2=\operatorname{Im}\tau.
\end{equation}
With the standard torus CFT normalization,
\begin{equation}
 \det{}'\Delta_{\can}=\tau_2|\eta(\tau)|^4,
 \qquad
 \boxed{Z_X^{\can,(1)}(\tau)
 =\frac{1}{\sqrt{\tau_2}\,|\eta(\tau)|^2}.}
 \label{eq:genus1-canonical}
\end{equation}
For $q=e^{2\pi i\tau}$, raw plumbing gives
\begin{equation}
 Z_X^{\pl,(1)}
 =\frac{1}{\sqrt{\tau_2}}
   \prod_{n=1}^{\infty}|1-q^n|^{-2}.
\end{equation}
Using $\eta(\tau)=q^{1/24}\prod_{n\geq1}(1-q^n)$,
\begin{equation}
 \boxed{
 W_1^{\pl/\can}=|q|^{1/12},
 \qquad
 \cA_1=\frac{1}{12}\log|q|.}
 \label{eq:genus1-weyl}
\end{equation}
Thus $Z_c^{\can}=Z_c^{\pl}|q|^{-c/12}$.  If the sewing propagator already
contains $q^{L_0-c/24}\bar q^{\bar L_0-c/24}$, this genus-one transition is
instead unity.

\section{The bielliptic conformal frame}
\label{sec:bielliptic-frame}

The analogue of the four-point pillow starts from a torus
\begin{equation}
 E_\tau=\mathbb C/(\mathbb Z+\tau\mathbb Z)
\end{equation}
with two insertions at $z_1,z_2$.  A connected double cover
$\pi:\Sigma_2\to E_\tau$ branched at these points has genus two by
Riemann--Hurwitz.  Its deck involution $\iota$ has two fixed points, and the
quotient $\Sigma_2/\langle\iota\rangle$ is the original torus with two
angle-$\pi$ corners.  Thus this is a branched conformal frame, rather than an
ordinary one-to-one change of coordinate.

Let $E(z,w)$ be the prime form of $E_\tau$; in the flat coordinate we use
\begin{equation}
 E(z,w)=\frac{\theta_1(z-w\mid\tau)}{\theta_1'(0\mid\tau)}.
\end{equation}
Choose $a\in E_\tau$ such that
\begin{equation}
 2a=z_1+z_2
 \qquad\text{in the group law of }E_\tau.
 \label{eq:half-branch-divisor}
\end{equation}
There are four choices, differing by $E_\tau[2]$; they encode the four
discrete $\mathbb Z_2$ monodromy sectors.  If
$z_1+z_2-2a=m+n\tau$, define
\begin{equation}
 f_a(z)=e^{-2\pi i n z}
 \frac{E(z,z_1)E(z,z_2)}{E(z,a)^2}.
 \label{eq:elliptic-cover-function}
\end{equation}
The exponential makes $f_a$ doubly periodic.  Its divisor is
$z_1+z_2-2a$, so the curve
\begin{equation}
 \boxed{\Sigma_{2,a}:\qquad \xi^2=f_a(z)}
 \label{eq:bielliptic-cover}
\end{equation}
is branched only at $z_1,z_2$; the double pole at $a$ is unramified.

The invariant and anti-invariant holomorphic one-forms are
\begin{equation}
 \omega_+=\pi^*\dd z,
 \qquad
 \omega_-=\pi^*\!\left(\frac{\dd z}{\xi}\right),
 \qquad
 \iota^*\omega_\pm=\pm\omega_\pm .
 \label{eq:eigen-one-forms}
\end{equation}
In the torus coordinate, the cut differential is
\begin{equation}
 \boxed{
 \omega_-(z)=\mathcal N_a e^{\pi i n z}
 \frac{\theta_1(z-a\mid\tau)}
 {\sqrt{\theta_1(z-z_1\mid\tau)\theta_1(z-z_2\mid\tau)}}\,\dd z .}
 \label{eq:cut-differential}
\end{equation}
For $z_1=0,z_2=u$, one may take
$a=u/2+\epsilon$ with $\epsilon\in E_\tau[2]$.  Formula
\eqref{eq:cut-differential} is the torus replacement of
$\dd z/\sqrt{z(z-1)(z-w)}$.  The theta-function numerator is required by
double periodicity; it also supplies the zero divisor of the genus-two
one-form.

Regularity is transparent locally.  At a branch point, write
$z-z_i=t^2$.  Then $\xi\sim t$ and
\begin{equation}
 \omega_-\sim\frac{\dd z}{t}=2\,\dd t.
\end{equation}
At $z=a$, one has $\xi\sim(z-a)^{-1}$ and hence
$\omega_-\sim(z-a)\dd z$.  The two lifts of $a$ therefore give the two
zeros required by $\deg K_{\Sigma_2}=2$.

Choose anti-invariant cycles $A_-,B_-$ and set
\begin{equation}
 \widehat\omega_-=
 \frac{\omega_-}{\oint_{A_-}\omega_-},
 \qquad
 \tau_{\mathrm P}(\tau,z_2-z_1)
 =\oint_{B_-}\widehat\omega_- .
 \label{eq:prym-period}
\end{equation}
The two eigenperiods are the quotient modulus $\tau$ and the Prym modulus
$\tau_{\mathrm P}$.  In a common involution-adapted convention,
\begin{equation}
 \Omega=\frac12
 \begin{pmatrix}
  \tau+\tau_{\mathrm P}&\tau-\tau_{\mathrm P}\\
  \tau-\tau_{\mathrm P}&\tau+\tau_{\mathrm P}
 \end{pmatrix}.
 \label{eq:bielliptic-period-matrix}
\end{equation}
Factors of two can move between \eqref{eq:prym-period} and
\eqref{eq:bielliptic-period-matrix} under a different integral homology
convention; invariantly, the two eigenperiods are $\tau$ and
$\tau_{\mathrm P}$.

Now present the genus-two surface as two pairs of pants joined by three
plumbing tubes $e_1,e_2,e_3$, with parameters $q_1,q_2,q_3$.  The
bielliptic involution exchanges $e_1,e_2$ and preserves $e_3$, so equivariant
sewing imposes
\begin{equation}
 \boxed{(q_1,q_2,q_3)=(x,x,y).}
 \label{eq:symmetric-plumbing-locus}
\end{equation}
The quotient graph is the two-edge necklace graph for a torus two-point
function: the orbit $\{e_1,e_2\}$ gives one quotient tube, the invariant
edge gives the other, and the two fixed points give the external insertions.
Equality of $q_1,q_2$ constrains the moduli, not the two channel labels; the
involution exchanges those labels, and an invariant block is obtained by
the corresponding projection.

Equations \eqref{eq:prym-period} and \eqref{eq:symmetric-plumbing-locus}
give an analytic definition of the desired change of variables:
\begin{equation}
 \tau_+(x,y)=\tau,
 \qquad
 \tau_-(x,y)=\tau_{\mathrm P}(\tau,u).
 \label{eq:period-matching-map}
\end{equation}
Equivalently, the two standard necklace multipliers are analytic functions
$p_1(x,y),p_2(x,y)$.  They need not equal $x,y$ literally: the quotient of
an invariant collar can introduce a square or square root, fixed by the
local plumbing coordinates.

Finally, the Abel--Jacobi metric diagonalizes in the eigenbasis as
\begin{equation}
 \dd s_{\can}^2\ \propto\
 \frac{|\omega_+|^2}{\operatorname{Im}\tau}
 +\frac{|\widehat\omega_-|^2}{\operatorname{Im}\tau_{\mathrm P}},
 \label{eq:bielliptic-canonical-metric}
\end{equation}
with the overall factor fixed by \eqref{eq:canonical-metric}.  Although
$\omega_-$ changes sign across the cut, $|\omega_-|^2$ is single-valued on
the quotient.  Near a fixed point,
\begin{equation}
 |\omega_-|^2\sim\frac{|\dd z|^2}{|z-z_i|},
\end{equation}
which is the angle-$\pi$ corner.  Thus the prime-form construction supplies
the conformal frame relating the torus two-point geometry to
$\Sigma_2/\mathbb Z_2$, as well as the explicit canonical metric needed for
the Weyl-anomaly quotient.

\section{Genus two}

Let $Y=\operatorname{Im}\Omega$ and use the raw Igusa product
\begin{equation}
 \chi_{10}^{\mathrm{raw}}(\Omega)
 =\prod_{\delta\ \mathrm{even}}
   \theta[\delta](0|\Omega)^2.
 \label{eq:chi10}
\end{equation}
Define the Petersson-norm combination
\begin{equation}
 F_2(\Omega)
 =(\det Y)^{5/2}
   \prod_{\delta\ \mathrm{even}}|\theta[\delta](0|\Omega)|
 =(\det Y)^{5/2}|\chi_{10}^{\mathrm{raw}}(\Omega)|^{1/2}.
 \label{eq:F2}
\end{equation}
For the Abel--Jacobi metric, Klein--Kokotov--Korotkin give
\begin{equation}
 \det{}'\Delta_{\can}=C_2 F_2(\Omega)^{1/3},
\end{equation}
where $C_2$ is independent of the moduli.  Hence
\begin{equation}
 \boxed{
 Z_X^{\can,(2)}(\Omega)
 =C_2^{-1/2}(\det Y)^{-5/12}
  |\chi_{10}^{\mathrm{raw}}(\Omega)|^{-1/12}.}
 \label{eq:genus2-canonical}
\end{equation}

Let $\Gamma_a$ be the rank-two Schottky group obtained from plumbing channel
$a$.  If $k_\gamma$ is the attracting multiplier of a primitive class, set
\begin{equation}
 \cP_a(\mathbf q)
 =\prod_{[\gamma]\in\mathcal P_{\mathrm{prim}}}
   \prod_{n=1}^{\infty}(1-k_\gamma(\mathbf q)^n)^{-1},
 \label{eq:schottky-product}
\end{equation}
with one representative of each inverse-paired primitive class.  The exact
raw plumbing scalar is
\begin{equation}
 \boxed{
 Z_X^{\pl,a}(\mathbf q)
 =(\det Y_a)^{-1/2}|\cP_a(\mathbf q)|^2.}
 \label{eq:genus2-plumbing}
\end{equation}
The factor $(\det Y_a)^{-1/2}$ is the Gaussian integral over the two handle
momenta; it is distinct from the constant target-space zero mode.

Dividing \eqref{eq:genus2-plumbing} by
\eqref{eq:genus2-canonical} gives the analytic Weyl factor
\begin{equation}
 \boxed{
 W_a^{\pl/\can}
 =C_2^{1/2}(\det Y_a)^{-1/12}
  |\chi_{10}^{\mathrm{raw}}(\Omega_a)|^{1/12}
  |\cP_a(\mathbf q)|^2.}
 \label{eq:genus2-weyl}
\end{equation}
Equivalently,
\begin{equation}
 \boxed{
 \cA_a
 =2\operatorname{Re}\log\cP_a
 -\frac{1}{12}\log\det Y_a
 +\frac{1}{12}\log|\chi_{10}^{\mathrm{raw}}(\Omega_a)|
 +\frac{1}{2}\log C_2.}
 \label{eq:genus2-log-weyl}
\end{equation}
Equations \eqref{eq:genus2-weyl}--\eqref{eq:genus2-log-weyl} are the desired
analytic answer.  The period map $\Omega_a(\mathbf q)$ and the multipliers
$k_\gamma(\mathbf q)$ have convergent Schottky expansions throughout the
plumbing domain.

\paragraph{Relative channel comparison.}
If channels $a$ and $b$ represent the same surface in the same symplectic
marking, then the canonical determinant and $C_2$ cancel:
\begin{equation}
 \frac{W_a^{\pl/\can}}{W_b^{\pl/\can}}
 =\frac{Z_X^{\pl,a}}{Z_X^{\pl,b}}.
 \label{eq:relative-channel}
\end{equation}
Period matrices recorded in different $\operatorname{Sp}(4,\mathbb Z)$
markings must first be transported to a common marking, together with the
corresponding automorphy factors.

\paragraph{Type 0B application.}
One real $N=1$ superfield has $c=3/2$, so its Weyl factor is
$(W_a^{\pl/\can})^{3/2}$.  Nine such superfields and the $b=1$
super-Liouville theory both have $c=27/2$; their common Weyl factor therefore
cancels in a same-frame quotient.

\section*{References}

\begin{enumerate}
 \item C.~Klein, A.~Kokotov, and D.~Korotkin,
 \emph{Extremal properties of the determinant of the Laplacian in the
 Bergman metric on the moduli space of genus two Riemann surfaces},
 \href{https://arxiv.org/abs/math/0511217}{arXiv:math/0511217}.
 \item G.~Mason and M.~P.~Tuite,
 \emph{The genus two partition function for free bosonic and lattice vertex
 operator algebras},
 \href{https://arxiv.org/abs/0712.0628}{arXiv:0712.0628}.
\end{enumerate}

\end{document}
